% Figure captions (greyscale figures; numbers come from numbers.tex)
\newcommand{\capMean}{(A) Mean body-weight-normalised vertical force (line) $\pm$ 1 SD (shaded) of the \nTrials{} analysed countermovement jumps. The take-off phase (movement onset to take-off) and the landing phase (touch-down to the lowest point of the centre of mass, COM) were each time-normalised to 0--100\% and concatenated; flight is omitted. The dotted line marks body weight. (B) Variance explained by each of the first ten principal components (bars) and cumulatively (line); the dashed line marks the four retained components.}
\newcommand{\capPCs}{Interpretation of the four retained principal components (main analysis). Each panel shows the mean curve (thin grey line) and the mean curve plus (solid black line) or minus (dashed black line) two standard deviations of the component score. Values below zero in the reconstructions are an artefact of adding a linear component to the mean and do not correspond to observed trials. The vertical line separates the take-off and landing phases.}
\newcommand{\capBand}{Trials matched on jump height do not share a common force--time curve. (A) Curves of the \jhBandN{} trials between the 45th and 55th percentiles of jump height (\jhBandMean{} $\pm$ \jhBandSD{}\,cm), shaded from light grey (low) to black (high) by their PC1 score. (B) Ratio of the standard deviation of each component score among trials matched on jump height (black circles) or on relative concentric peak force (grey triangles) to that in the full sample, with athlete-level bootstrap 95\% confidence intervals. A ratio of 1 indicates that matching on the metric leaves curve-shape variation unchanged.}
\newcommand{\capRtwo}{Proportion of the variance of each component score explained by the four vendor metrics, age and body mass, under four evaluation schemes (bars from black to white: in-sample, athlete-grouped cross-validation, one trial per athlete, athlete means), for the main landing definition (left) and the 1-s landing window (right). CV, cross-validation with athletes allocated to folds.}
\newcommand{\capWindow}{Sensitivity analysis with the landing phase ending 1\,s after touch-down: mean curve (thin grey line) and the mean curve plus (solid black line) or minus (dashed black line) two standard deviations of each retained component. Under this definition PC3 is located almost entirely in the take-off phase (\ldsSPCCtakeoff\% of its squared loading) and describes the duration and amplitude of the propulsive force.}
